\section{Discussion}
\label{sec:discussion}

\textbf{Mitigation effectiveness is not uniform.} Our results demonstrate that safety system prompts significantly improve mitigation robustness across all tested models ($p < 0.01$), but their effect on other dimensions---calibration, evidence grounding, and refusal consistency---varies substantially. This suggests that system-prompt-based mitigation is a necessary but insufficient layer: it reliably shifts behavior on the dimension it most directly targets while having unpredictable side effects on subtler dimensions.

\textbf{Domain-specific patterns reveal structural insights.} The domain-family analysis (\Cref{sec:domain_analysis}) shows that mitigation-sensitive query handling---items testing boundary maintenance under reframing---benefits most from system prompts ($\Delta$ up to $+1.25$). Conversely, ambiguity recognition items show degradation for four of six models (GPT-4o, DeepSeek-V3, Gemini-2.5-Pro, Qwen3-32B), suggesting that safety prompts may suppress appropriate uncertainty expression. This motivates domain-aware mitigation strategies rather than uniform system prompts.

\textbf{Strong baselines reduce marginal mitigation benefit.} GPT-4o's near-zero overall delta ($+0.02$, 95\% CI $[-0.13, +0.18]$) does not indicate that mitigation is ineffective---rather, that its base safety training already captures most behavior the mitigation prompt would induce. The pattern of lower baseline yielding larger marginal gain ($r = -0.67$ across models) holds regardless of whether the model is open-weight or closed-source. The practical implication is that system-prompt mitigation is most valuable for models with weaker built-in safety training.

\textbf{Calibration--safety trade-offs deserve explicit monitoring.} The significant calibration decrease for GPT-4o under mitigation ($\Delta = -0.38$, $p = 0.036$, $d = -0.49$) is not a measurement artifact---it reflects a genuine tension between safety-oriented prompting and epistemic calibration. Safety prompts that encourage caution may lead models to hedge more than evidence warrants. Benchmark designers should monitor this trade-off explicitly rather than relying on aggregate safety scores.

\textbf{Policy inconsistency as a persistent failure mode.} The \texttt{policy\_inconsistency} error tag persists across all models even after mitigation, indicating that models apply safety thresholds inconsistently across semantically equivalent prompts. This is a structural limitation of current architectures: system prompts shift mean behavior but do not eliminate sensitivity to surface-level variation in prompt phrasing.

\subsection{Limitations}

\textbf{Scale.} The public benchmark contains 24 items. While this provides balanced coverage across 6 domain families and 6 reasoning types, it limits statistical power---the Friedman tests for cross-model differences are non-significant ($p > 0.38$), likely reflecting insufficient power rather than true model equivalence. The restricted layer (not released) is larger but cannot be used for open benchmarking.

\textbf{Scoring methodology.} The primary scoring protocol uses pattern-based indicators. While validated against an independent LLM-as-judge (\Cref{sec:llm_judge}), neither approach replicates expert human evaluation with domain-specific CBRN knowledge. Future work should validate against independent expert ratings.

\textbf{Model selection.} We evaluate six models at a single point in time. Model behavior changes across versions, and our results reflect a snapshot rather than a longitudinal trend. The framework supports longitudinal tracking, but this paper reports only the initial evaluation round.

\textbf{Single mitigation strategy.} We test only system-prompt-based mitigation. Other strategies (fine-tuning, content filtering, retrieval-augmented safety) may produce different patterns and should be evaluated using the same framework.
